\appendix
\section{Reproduction}
\label{app:repro}

The companion repository~\cite{rhoblackwell}, \url{https://github.com/aburan28/rho-blackwell}, carries the
ECC2K-130 client (CUDA and CPU), the NIST P-256/P-384 rho kernel, the
paper source, the benchmark data behind every figure, the evidence
records of Appendix~\ref{app:evidence}, and a Dockerfile.

\paragraph{Status.} The CPU client, its validation suite, a
planted-logarithm solve and every mode of the image's entry point have
been run from the repository's copy of the source; the Docker image
itself has not yet been built, and its GPU modes have not been run. The
commands below describe the intended workflow.

\paragraph{Docker.} The image compiles the client with CUDA~13.3 for
Ampere, Ada, Hopper and both Blackwell families, with the audited
RTX~PRO~6000 preset (the $\sigma$-walk, batch 16, 256-thread blocks,
native \texttt{clmad}, compact tiled state and shared Frobenius masks),
and runs the client's validation suite during the build. By default a
container walks ECC2K-130 with the campaign's parameters (cutoff 32, a
$2^{32}$ iteration guard, checkpoints every ten minutes) on the GPU if one
is visible and on the CPU otherwise, writing 32-byte distinguished-point
records and a resumable checkpoint to its data volume:
{\small\begin{verbatim}
docker build -t rho-blackwell -f docker/Dockerfile .
docker run --gpus all -v rho-data:/data rho-blackwell   # walk ECC2K-130
docker run --gpus all rho-blackwell bench                # throughput only
docker run rho-blackwell test                             # validation suite
docker run rho-blackwell planted                          # small planted logs
docker run --gpus all rho-blackwell ecc2k95               # ECC2K-95, answer known
docker run --gpus all rho-blackwell nist                  # P-256/P-384 rho bench
\end{verbatim}}
A container without \texttt{RUN\_ID} picks a random run identifier from
a range the campaign never used (20000--65535), keeps it with its data
and warns that it did so. Two containers with the same identifier walk
the same trails, and with random choices that happens with probability
about $0.1\%$ for 10 containers, $10\%$ for 100 and $63\%$ for 300, so any
run of more than one container must assign each a distinct
\texttt{RUN\_ID}. A collision is
found by merging the corpora of all containers that walked with the
same cutoff; a single container is unlikely to find one alone.

\paragraph{Building without Docker.} The client needs a C++17 compiler
and, for the GPU build, CUDA~13.3 or newer (the native carry-less product
requires PTX~9.3). Its make targets build the CPU client, the GPU client
with any preset of this paper, and the test suites; the repository's
README lists the targets for each measured build, including the fused
and nibble-table $\sigma$-walks and the B200 and RTX~PRO~4500 presets.

\paragraph{Checking a number.} Every measured number cites an evidence
record [E$n$]. The repository's \texttt{evidence/} directory holds each
record under its number, with the notes, data or source that report the
measurement, so a reader can check a number without the development
history; \texttt{evidence/map.json} names the original file of each.
